\section{Decoder structure and polynomial implementation}
\label{sec:decoder-structure}
The answer space contains $2^M$ maps, where $M=|\calC|$. Its coordinates are
indexed by disjoint observational classes, which gives a factorization that is
not available for an arbitrary multiple-answer identification problem.

\subsection{Computing the coefficient without enumerating answers}
For a class $C$ and a candidate output $\ell\in\{0,1\}$, put
\[
r_i(C,\ell)=\min_{j\in C:\theta_j\ne\ell}\KL(P_i\Vert P_j),
\qquad
\mathcal L_i(C)=
\begin{cases}
\{\theta_i\},&C=C_i,\\
\{0,1\},&C\ne C_i.
\end{cases}
\]
The empty minimum is $+\infty$.
\begin{theorem}[Factorized decoder coefficient]
\label{thm:factorization}
Under Assumption~\ref{ass:repair},
\begin{equation}
D_i=\min_{C\in\calC}\ \max_{\ell\in\mathcal L_i(C)}r_i(C,\ell).
\label{eq:factorization}
\end{equation}
An optimal decoder can be obtained by selecting, independently in each class,
any allowed $\ell$ with $r_i(C,\ell)\ge D_i$. Choosing the smallest such $\ell$
in every coordinate gives the lexicographically first global optimum.
Given the pairwise KL matrix, all $D_i$ and a corresponding decoder for each
environment require $O(K^2+KM)$ arithmetic operations.
\end{theorem}
\begin{proof}
For every decoder valid in $i$,
\[
d_i(a)=\min_{C\in\calC}r_i(C,a(C)).
\]
Its coordinate in class $C$ belongs to $\mathcal L_i(C)$. Thus
$d_i(a)\le\min_C\max_{\ell\in\mathcal L_i(C)}r_i(C,\ell)$.
Conversely choose a maximizing label in each class. All choices satisfy
validity in $i$, and their minimum equals the displayed upper bound.
This proves equality. Once $D_i$ is known, a valid decoder is globally optimal
exactly when every coordinate meets this threshold. Coordinatewise smallest
feasible choices therefore yield the first lexicographic optimum; requiring
every coordinate to be a local maximizer would unnecessarily discard some
global optima.

For one row of the KL matrix, scan the $K$ environments and keep, in each
class, the smallest divergence for each of the two labels. These give every
$r_i(C,\ell)$ in $O(K+M)$ operations. Scan the classes again to take the
bottleneck and construct a decoder. Repeat for all $i$. Computing the KL
matrix itself uses $O(K^2|\mathcal X_{\rm pre}|)$ operations for a finite
pre-repair alphabet, before accounting for numerical precision.
\end{proof}

For a class containing only one old label, choosing that label incurs no
invalidating alternative and contributes $+\infty$ to the bottleneck.
Mixed-label classes are therefore the only effective constraints. In class
$C_i$ the output is fixed to $\theta_i$; other classes may be assigned whichever
label leaves the easier set of alternatives to exclude. This is the mechanism
behind same-label exclusion in Section~\ref{sec:example}.

\subsection{A surviving-set test for the likelihood-ratio race}
Let $s_j(n)=\sum_{t=1}^n\log P_j(X_t)$ and define, for a candidate witness $k$,
\[
S_k(n)=\{j:s_k(n)-s_j(n)<\beta_\delta\}.
\]
It includes $k$, since $\beta_\delta>0$.

\begin{theorem}[Polynomial eligibility test]
\label{thm:polynomial}
A decoder eligible with witness $k$ in \eqref{eq:race} exists if and only if
every nonempty intersection $S_k(n)\cap C$ has a single old label.
For such a witness, each nonempty intersection forces that label as the
coordinate $a(C)$; every empty intersection leaves the coordinate free.
Consequently, the lexicographically first eligible decoder can be found in
$O(K^2+KM)$ operations per observation, without enumerating $2^M$ maps.
For the separated race, eligibility is equivalent to all members of $S_k(n)$
having one common old label.
\end{theorem}
\begin{proof}
The likelihood condition requires every invalidating environment $j\in B(a)$
to lie outside $S_k(n)$. Equivalently, $a(C_j)=\theta_j$ for every surviving
$j$. One coordinate can satisfy these requirements precisely when the labels
within each surviving class intersection agree. Since $k$ survives, the
resulting decoder is automatically valid at its witness.

For each witness, assign the forced labels and zero to unconstrained
coordinates. This is the smallest eligible decoder for that witness. Taking
the smallest of these candidates, with the witness index as a secondary
tie-break, returns exactly the first decoder of an exhaustive lexicographic
race. A scan over the $K$ environments followed by the $M$ coordinates
suffices for each witness. Constant decoders satisfy all surviving constraints
exactly when their labels agree globally.
\end{proof}

This implementation realizes the same statistical rule analyzed below; it
does not change its threshold or error guarantee. An online controller stores
$K$ log-likelihood scores, the class and label maps, and a candidate decoder.
Its working storage can be $O(K+M)$ when witnesses are checked successively.
Floating-point equality at a likelihood threshold is still a numerical issue:
the exact rule uses the displayed strict surviving-set inequality, and the
implementation uses the complementary comparisons consistently.

The finite Bernoulli evaluator batches many possible count histories, so its
time and storage costs differ from those of a single online trajectory.
Its bit-coded batch implementation supports at most 60 classes; the scalar
polynomial implementation has no such encoding restriction.
